% =====================================================================
%  Soutenance : Analyse de la dynamique fractionnaire de la mobilité urbaine
%  Compilation : latexmk -pdf -xelatex slide.tex      (figures dans ./figures/)
% =====================================================================
\documentclass[aspectratio=169,11pt]{beamer}

\usepackage{fontspec}
\usepackage[french]{babel}
\usepackage{amsmath,amssymb}
\usepackage{booktabs,array}
\usepackage{graphicx}
\graphicspath{{figures/}}
\usepackage{tikz}
\usetikzlibrary{arrows.meta,positioning}   % <-- AJOUT : positioning (corrige l'erreur PGF)

\usetheme{Madrid}
\definecolor{accent}{HTML}{1F4E79}
\setbeamercolor{structure}{fg=accent}
\setbeamercolor{alerted text}{fg=red!70!black}
\setbeamertemplate{navigation symbols}{}
\setbeamertemplate{footline}{\hfill\small\insertframenumber\hspace{4mm}\vspace{2mm}}
\setbeamertemplate{itemize items}[circle]

\newcommand{\code}[1]{\texttt{#1}}
% \fig[hauteur max en fraction de \textheight]{fichier}
\newcommand{\fig}[2][0.72]{\includegraphics[width=\linewidth,height=#1\textheight,keepaspectratio]{#2}}

\title[Dynamique fractionnaire de la mobilité]{Analyse de la dynamique fractionnaire\\de la mobilité urbaine}
\subtitle{Exposant de Hurst, temps de séjour et lois de puissance\\sur des trajectoires GPS (GeoLife)}
\author{Votre nom}            % <- à compléter
\institute{Master 2 \,---\, 2025--2026}
\date{Soutenance}

\begin{document}

% ---------------------------------------------------------------------
\begin{frame}[plain]
\titlepage
\end{frame}

\begin{frame}{Plan}
\tableofcontents
\end{frame}

% =====================================================================
\section{Contexte et problématique}
% =====================================================================

\begin{frame}{Contexte : la mobilité humaine vue par GPS}
\begin{itemize}
  \item Récepteurs GPS embarqués : suivi de la mobilité à l'échelle de quelques secondes
  \item Résultats classiques : \alert{lois d'échelle} des longueurs de déplacement, forte régularité des trajectoires individuelles\\
        {\footnotesize (Brockmann 2006 ; González 2008 ; Rhee 2011)}
  \item Notre angle : la \alert{mémoire}
\end{itemize}
\begin{block}{Question}
L'amplitude d'un déplacement influence-t-elle les suivants, au-delà des corrélations de court terme ?
\end{block}
\medskip
Deux axes dans ce travail :\ \textbf{(1)} l'exposant de Hurst $H$ ;\ \textbf{(2)} les lois de puissance des temps de séjour et des distances entre arrêts.
\end{frame}

\begin{frame}{L'exposant de Hurst $H$}
\begin{columns}[T]
\column{0.50\textwidth}
\begin{tabular}{@{}lp{5.4cm}@{}}
\toprule
$H$ & Lecture \\
\midrule
$H<\tfrac12$ & antipersistant \\
$H=\tfrac12$ & sans mémoire (bruit blanc) \\
$H>\tfrac12$ & persistant : un grand déplacement tend à être suivi d'un grand \\
\bottomrule
\end{tabular}
\column{0.47\textwidth}
\begin{block}{Dépendance à longue portée (DLP)}
\[
\gamma(k)\sim C\,k^{2H-2},\qquad \tfrac12<H<1
\]
$\sum_k|\gamma(k)|=+\infty$ : les corrélations décroissent si lentement que leur somme diverge.\\[1mm]
Spectre : $f(\lambda)\sim|\lambda|^{1-2H}$ quand $\lambda\to0$.
\end{block}
\end{columns}
\medskip
Modèle de référence : le \alert{bruit gaussien fractionnaire} (fGn), qui sert à fabriquer des séries dont $H$ est connu.
\end{frame}

\begin{frame}{Positionnement}
\begin{itemize}
  \item Travail le plus proche : Wang, Han \& Wang (2014) estiment $H$ sur des séries de déplacements issues de GeoLife, construites à partir des \alert{déplacements entre points d'arrêt successifs}
  \item Notre série : les \alert{distances entre points GPS successifs}, $d_i$
  \item Cette différence soulève des difficultés méthodologiques qui forment le fil conducteur du mémoire :
  \begin{itemize}
    \item non-stationnarité
    \item changements de régime
    \item rythme d'enregistrement
  \end{itemize}
\end{itemize}
\end{frame}

\begin{frame}{Du cahier des charges à une hypothèse à tester}
\textbf{Demandé :} \og démontrer empiriquement \fg{} la DLP, et mettre en évidence des lois de puissance.

\medskip
\textbf{Trois obstacles à une démonstration directe :}
\begin{enumerate}
  \item \alert{La série} : $d_i=v_i\,\delta_i$ mélange vitesse et intervalle de temps ; elle n'est pas stationnaire
  \item \alert{Les estimateurs} : sensibles à la non-stationnarité et aux changements de régime, faciles à confondre avec de la mémoire longue (Diebold \& Inoue 2001)
  \item \alert{La circularité} : le filtre de fiabilité retient les segments où les méthodes s'accordent, donc conditionne sur $H$ lui-même
\end{enumerate}
\begin{block}{Reformulation}
$H_0$ : mémoire courte avec changements de régime \quad contre \quad $H_1$ : mémoire longue
\end{block}
\end{frame}

\begin{frame}{Question de recherche et périmètre}
\begin{block}{Question}
Dans quelle mesure peut-on, après prétraitement, attribuer à la mobilité la valeur de $H$ estimée, et que permet de conclure un test à mémoire courte ?
\end{block}
\footnotesize
\begin{tabular}{@{}p{8.2cm}l@{}}
\toprule
Composant & État \\
\midrule
Prétraitement, marquage, indicateurs de qualité & implémenté \\
Détection d'arrêts et temps de séjour & implémenté, limites connues \\
Estimation de $H$ (DFA, R/S, Whittle local) et filtre de fiabilité & implémenté \\
Validation synthétique, test de dépassement & implémenté \\
Lois de puissance (Clauset, alternatives) & méthode et résultats \\
Test du mécanisme ON/OFF & script écrit, à corriger \\
Segmentation par labels, séries à pas régulier, tests de Berkes et de Qu & proposé \\
Typologie des utilisateurs, cartographie des zones de mémoire & non commencé \\
\bottomrule
\end{tabular}
\end{frame}

% =====================================================================
\section{Données}
% =====================================================================

\begin{frame}{GeoLife en chiffres}
Microsoft Research Asia. Guide v1.2 : 178 utilisateurs, 17\,621 trajectoires, 23\,667\,828 points, 1\,251\,654~km, 48\,203~heures.
\bigskip
\begin{columns}[c]
\column{0.3\textwidth}\centering {\Huge\textbf{145}}\\[1mm] utilisateurs avec au moins\\ un segment retenu
\column{0.3\textwidth}\centering {\Huge\textbf{19\,150}}\\[1mm] segments retenus
\column{0.3\textwidth}\centering {\Huge\textbf{7\,853}}\\[1mm] segments \og gardés \fg{}\\ (41,0~\%)
\end{columns}
\bigskip
\begin{itemize}
  \item Aucun protocole de collecte : chacun enregistrait quand et comme il le voulait
  \item Dossier 182 : données personnelles de l'auteur (Sénégal, période récente) \(\Rightarrow\) \alert{exclu} de toute analyse agrégée
\end{itemize}
\end{frame}

\begin{frame}{Format des fichiers et mode d'enregistrement}
\begin{itemize}
  \item Un fichier \code{.plt} par trajectoire : 6 lignes d'en-tête, puis latitude, longitude, altitude (pieds, $-777$ si invalide), jour julien, date et heure (GMT)
  \item 91~\% des trajectoires sont denses : un point toutes les \alert{1 à 5~s}, ou tous les \alert{5 à 10~m}
  \item Ces deux modes ne mesurent pas la même chose : en mode \og distance \fg{}, $d_i$ est presque constant quelle que soit la vitesse \(\Rightarrow\) la série décrit l'appareil, pas le déplacement
  \item Heuristique de classement : mode \emph{temps}, \emph{distance} ou \emph{mixte}
  \item Labels de transport : 69 utilisateurs (saisis à la main), \alert{non utilisés} dans le pipeline actuel
\end{itemize}
\end{frame}

\begin{frame}{Audit de T-Drive : un second jeu de données écarté}
\begin{columns}[T]
\column{0.58\textwidth}
\fig[0.65]{graphe_6_tdrive.pdf}        % <-- MODIFIÉ : 0.70 -> 0.65
\column{0.40\textwidth}
\begin{itemize}
  \item $65{,}3$~\% des taxis dépassent 200~km/h, $48{,}6$~\% dépassent 1\,000~km/h
  \item Sauts de 1\,495~km sans déplacement physique (taxi 6275)
  \item Taille médiane des segments à $T_{\mathrm{gap}}=900$~s : \alert{10 points}
  \item Lois de puissance : $x_{\min}=12{,}7$~h, queue de $1{,}8$~\% des séjours (stationnements nocturnes)
\end{itemize}
\medskip
\alert{T-Drive mesure l'exploitation d'une flotte, pas la mobilité.}
\end{columns}
\end{frame}

% =====================================================================
\section{Méthodes}
% =====================================================================

\begin{frame}{Vue d'ensemble de la chaîne}
\centering
\resizebox{\linewidth}{!}{%
\begin{tikzpicture}[
  box/.style={draw=accent, rounded corners=2pt, fill=accent!8, align=center, font=\scriptsize, minimum height=9mm, inner sep=2pt},
  b1/.style={box, text width=2.05cm},
  b2/.style={box, text width=1.85cm},
  arr/.style={-{Stealth[length=1.6mm]}, thick, accent}]
\node[b1] (a) at (1.05,0)  {Fichiers \code{.plt}\\GeoLife};
\node[b1] (b) at (3.55,0)  {Nettoyage,\\tri, doublons};
\node[b1] (c) at (6.05,0)  {Segmentation\\(trou $>900$~s)};
\node[b1] (d) at (8.55,0)  {Marquage des\\points suspects};
\node[b1] (e) at (11.05,0) {Sélection\\($\ge300$ points)};
\draw[arr] (a)--(b); \draw[arr] (b)--(c); \draw[arr] (c)--(d); \draw[arr] (d)--(e);
\node[b2] (f) at (1.0,-2.2)  {Série $d_i$\\(distances)};
\node[b2] (g) at (3.3,-2.2)  {Estimation de $H$\\DFA, R/S, Whittle};
\node[b2] (h) at (5.6,-2.2)  {Filtre\\\og gardé \fg{}};
\node[b2] (i) at (8.3,-2.2)  {Détection d'arrêts\\200~m, 20~min};
\node[b2] (j) at (10.6,-2.2) {Séjours et\\distances};
\node[b2] (k) at (12.9,-2.2) {Clauset +\\4 alternatives};
\draw[arr] (f)--(g); \draw[arr] (g)--(h); \draw[arr] (i)--(j); \draw[arr] (j)--(k);
\draw[accent, thick] (e.south) -- ++(0,-0.55) coordinate (m);
\draw[arr] (m) -| (f.north);
\draw[arr] (m) -| (i.north);
\node[font=\tiny, below=1mm of g] {+ test de dépassement AR($p$)};
\node[font=\small\bfseries, accent] at (3.3,-3.6) {Branche $H$};
\node[font=\small\bfseries, accent] at (10.6,-3.6) {Branche lois de puissance};
\end{tikzpicture}}
\end{frame}

\begin{frame}{Prétraitement : quatre principes}
\begin{enumerate}
  \item \alert{Marquer avant de décider} : les points suspects sont signalés sans être retirés ; l'effet d'une coupure est mesuré à part
  \item \alert{Seuils relatifs au contexte du segment}, bornés par des planchers ou plafonds absolus
  \item \alert{Deux versions de chaque estimateur} : référence lisible et vectorisée, reliées par des tests de parité
  \item \alert{Traçabilité} : empreinte de la configuration sur chaque ligne de résultats, graines fixées, tableaux produits par script
\end{enumerate}
\medskip
Lecture : conversion en secondes UNIX (UTC), rejet des coordonnées hors bornes, tri par horodatage, suppression des doublons (premier conservé).
\end{frame}

\begin{frame}{Segmentation et calibration des seuils}
\begin{columns}[T]
\column{0.55\textwidth}
\fig[0.78]{graphe_5_heatmap.pdf}
\column{0.43\textwidth}
\begin{itemize}
  \item Un trou $>T_{\mathrm{gap}}=900$~s ouvre un nouveau segment
  \item \alert{Ni interpolation, ni raccord} : un saut artificiel fausse la DFA
  \item Distribution des $\Delta t$ : aucun pic naturel \(\Rightarrow\) seuil \alert{externe}
  \item Couple retenu $(900,\,300)$ : \textbf{91,3~\%} des points conservés ; médiane de \textbf{721 points} par segment
  \item Décision d'ingénierie, entourée de valeurs proches (89,6 à 600~s ; 92,2 à 1\,200~s)
\end{itemize}
\end{columns}
\end{frame}

\begin{frame}{Marquage des points suspects}
\begin{columns}[T]
\column{0.52\textwidth}
\footnotesize
\begin{tabular}{@{}lp{4.6cm}@{}}
\toprule
Composant & Seuil \\
\midrule
Saut de distance & $\min(30\,\tilde d,\,500)$ m \\
Vitesse maximale & $\min\bigl(200,\max(5,20\,\tilde v)\bigr)$ km/h \\
Variation brutale & $\max(20,\,5\,\mathrm{MAD}_v)$ km/h \\
Persistance (régime) & $\max(5,\mathrm{round}(0{,}01N))$ points \\
\bottomrule
\end{tabular}
\medskip
\normalsize
Candidat : Z-score modifié $>3{,}5$.
\column{0.46\textwidth}
\begin{block}{Classification}
\small
\begin{itemize}
  \item brutal et isolé $\Rightarrow$ \code{artefact}
  \item soutenu $\Rightarrow$ \code{regime}
  \item modéré et isolé $\Rightarrow$ non marqué
\end{itemize}
\end{block}
\footnotesize
Volumes : \code{suspect\_dist} 148\,203 (0,67~\%) ; \code{suspect\_vitesse} 273\,499 (1,25~\%).\\[1mm]
\alert{Faiblesse} : dans un segment piéton (médiane 3~km/h), le seuil de vitesse tombe vers 60~km/h ; un trajet motorisé soutenu risque d'être classé en artefact.
\end{columns}
\end{frame}

\begin{frame}{La série analysée : les distances $d_i$}
\begin{block}{Ce que mesure $d_i$}
$d_i=v_i\,\delta_i$ : la distance entre deux points GPS successifs mélange \alert{vitesse} et \alert{intervalle de temps}.
\end{block}
\begin{itemize}
  \item Imposée par le cahier des charges ; mesure l'\alert{intensité} du déplacement, pas sa direction
  \item Mode d'enregistrement \og distance \fg{} : la série décrit l'appareil
  \item Contrôle prévu : analyse de la série des vitesses $v_i$ (option du code)
  \item Stationnarité : test KS entre les deux moitiés ; $p<10^{-20}$ sur la quasi-totalité des segments gardés \(\Rightarrow\) indicateur peu discriminant
\end{itemize}
\end{frame}

\begin{frame}{Trois estimateurs de $H$}
\footnotesize
\begin{columns}[T]
\column{0.32\textwidth}
\begin{block}{DFA}
Profil $Y(i)=\sum_{k\le i}(x_k-\bar x)$, fenêtres disjointes, droite retirée dans chacune.\\[1mm]
$F(s)\sim s^{\alpha}$ : $\alpha=H$ (fGn), $\alpha=H+1$ (fBm).\\[1mm]
20 échelles de $s_{\min}=4$ à $N/4$.
\end{block}
\column{0.32\textwidth}
\begin{block}{R/S}
Étendue $R(s)$ de la somme cumulée sur l'écart-type $S(s)$ :\\
$\mathbb E[R/S]\sim c\,s^{H}$.\\[1mm]
Pente \alert{brute} et pente \alert{corrigée} (Anis--Lloyd).\\[1mm]
Blocs de 10 à $N/2$.
\end{block}
\column{0.32\textwidth}
\begin{block}{Whittle local}
Basses fréquences seulement ($m=N^{0{,}65}$) ; minimisation de $R(H)$ sur la grille $H_k=k/1001$.\\[1mm]
Indicateur \code{whittle\_a\_la\_borne} : minimum sur une borne $\Rightarrow$ ce n'est pas une estimation.
\end{block}
\end{columns}
\end{frame}

\begin{frame}{Validation sur séries de $H$ connu (fGn)}
\begin{columns}[T]
\column{0.60\textwidth}
\fig[0.78]{graphe_1_validation.pdf}
\column{0.38\textwidth}
\small
\begin{itemize}
  \item 200 réalisations par cellule ; 2\,400 séries ($4$ valeurs de $H_0$, $3$ longueurs)
  \item DFA et Whittle : biais proche de zéro ; RMSE $\le0{,}067$ (DFA) et $\le0{,}095$ (Whittle) à $N=300$
  \item R/S corrigé : fiable pour $H\le0{,}5$, \alert{biais de $-0{,}149$} à $H_0=0{,}9$, $N=300$ : la correction sur-corrige les séries persistantes
  \item \(\Rightarrow\) ne pas prendre le R/S comme référence indépendante pour valider un $H$ élevé
\end{itemize}
\end{columns}
\end{frame}

\begin{frame}{Test de dépassement d'un modèle à mémoire courte}
\begin{enumerate}
  \item Transformer la série en loi normale à partir de ses rangs
  \item Ajuster un AR($p$) par Yule--Walker, $p\le p_{\max}$ choisi par l'AIC
  \item Simuler 99 séries AR($p$) de même longueur
  \item Leur redonner la distribution de valeurs de la série observée (échange de rangs)
\end{enumerate}
\medskip
$p$-valeur unilatérale : $p=\dfrac{1+\#\{h_{\mathrm{surr}}\ge h_{\mathrm{obs}}\}}{B+1}$
\begin{alertblock}{Lecture}
Le rejet dit : \og la dépendance n'est \emph{pas expliquée} par un AR($p$) \fg{}. Ce n'est \alert{pas} une preuve de mémoire longue : une alternance de régimes produit le même rejet.
\end{alertblock}
\end{frame}

\begin{frame}{Le filtre de fiabilité : critère \og gardé \fg{}}
Un segment est \alert{gardé} si et seulement si :
\begin{enumerate}
  \item les trois méthodes ont produit une valeur
  \item $0<\hat H<1$ pour chacune des trois méthodes
  \item le minimum de $R(H)$ du Whittle n'est pas sur une borne de la grille
  \item la divergence (max $-$ min des trois estimations) est au plus de $0{,}15$
\end{enumerate}
\medskip
\begin{itemize}
  \item Drapeau \code{suspect} : divergence $>0{,}15$, ou $\hat H_{\mathrm{DFA}}>1$, ou $\hat H_{\mathrm{RS}}>1$, ou $\hat H_{\mathrm{Whittle}}>0{,}999$
  \item Résultat : \textbf{7\,853} segments gardés sur 19\,150 (41,0~\%) ; 97,1~\% d'entre eux ont les trois méthodes $>0{,}5$
\end{itemize}
\end{frame}

\begin{frame}{Limite principale : la circularité du filtre}
\begin{columns}[T]
\column{0.50\textwidth}
\fig[0.78]{graphe_3_circularite.pdf}
\column{0.48\textwidth}
\begin{itemize}
  \item Le filtre porte sur le résultat : la divergence utilise les $\hat H$ eux-mêmes
  \item $|\hat H_{\mathrm{DFA}}-\hat H_{\mathrm{RS}}|$ moyen :\\ $0{,}079\to0{,}023$ (par utilisateur)\\ $0{,}152\to0{,}058$ (par segment)
  \item Réduction \alert{acquise par construction}, pas une confirmation indépendante de l'accord des méthodes
\end{itemize}
\begin{alertblock}{À retenir}
$\hat H_{\mathrm{gardés}}$ est un $\hat H$ \alert{conditionnel au filtre}, pas un $\hat H$ de population.
\end{alertblock}
\end{columns}
\end{frame}

\begin{frame}{Arrêts et temps de séjour}
\begin{itemize}
  \item \alert{Fenêtre ancrée} (Li et al.) : pour chaque point $p_i$, premier $j$ tel que $\mathrm{dist}(p_i,p_j)>R_{\mathrm{arr}}$ ; si $t_j-t_i\ge T_{\min}$, arrêt localisé au centroïde
  \item Réglage : $R_{\mathrm{arr}}=200$~m, $T_{\min}=1\,200$~s (20 min), école de Zheng et al. ; sensibilité sur 4 couples
  \item \alert{Censure} : durées biaisées \emph{étiquetées}, jamais écartées silencieusement ; par défaut, ajustement sur les arrêts internes non censurés
  \item Deux variables de distance :
  \begin{itemize}
    \item \code{distance} : même segment, 1\,283 paires (environ 87~\% des paires écartées)
    \item \code{distance\_a\_cheval} : inclut les trous d'enregistrement, 10\,216 paires
  \end{itemize}
  \item Limites : détection sur les segments déjà retenus ; fragmentation d'un long séjour par le bruit GPS ; durée légèrement surestimée
\end{itemize}
\end{frame}

\begin{frame}{Lois de puissance : protocole}
\begin{columns}[T]
\column{0.50\textwidth}
\begin{block}{Ajustement de Clauset}
\small
\begin{itemize}
  \item Maximum de vraisemblance de $\alpha$
  \item $x_{\min}$ par distance KS (200 candidats, queue $\ge50$)
  \item Adéquation : bootstrap semi-paramétrique, $B=1\,000$
  \item Rejet si $p\le0{,}1$
\end{itemize}
\end{block}
\column{0.48\textwidth}
\begin{block}{Alternatives}
\small
\begin{itemize}
  \item Exponentielle, log-normale, Weibull, loi \alert{tronquée} $x^{-\alpha}e^{-\lambda x}$
  \item Vuong (non emboîtés), $\chi^2$ (tronquée vs pure)
  \item Poids d'Akaike corrigés (AICc)
\end{itemize}
\end{block}
\end{columns}
\begin{alertblock}{Adéquation de la tronquée}
Poids AICc = meilleur modèle \emph{du catalogue}, pas bon ajustement. Test bootstrap ($B=200$, $x_{\min}$ fixé) : moins exigeant que Clauset, donc $p>0{,}1$ ne prouve rien, tandis que $p\le0{,}1$ est un rejet net.
\end{alertblock}
\end{frame}

% =====================================================================
\section{Résultats sur l'exposant de Hurst}
% =====================================================================

\begin{frame}{Sept hypothèses, sept verdicts}
\footnotesize
\begin{tabular}{@{}l p{4.6cm} p{6.4cm}@{}}
\toprule
 & Énoncé & Verdict \\
\midrule
H1 & Estimateurs fiables sur fGn & DFA et Whittle confirmés ; R/S infirmé pour $H\ge0{,}7$ \\
H2 & Un régime fait rejeter le nul AR($p$) & atteint pour $p_{\max}=1$ (100~\% à $N=1000$), pas pour $p_{\max}=5$ (19~\%) \\
H3 & Corrélation $\rho_1$--$H$ & confirmée ($\rho_S=0{,}859$) \\
H4 & Un AR(1) reproduit les taux du corpus & infirmée (corpus 20,3~\% ; AR(1) : 2 à 6~\%) \\
H5 & Mobile / non mobile à $\rho_1$ fixé & confirmée mais fragile ($\Delta=+0{,}057$) \\
H6 & Test nul calibré à $\rho_1=0$ & confirmée (5,25~\% de faux positifs) \\
H7 & Corpus invalide pour $H$ & infirmée (82,9~\% de segments valides) \\
\bottomrule
\end{tabular}
\medskip
\normalsize
Le taux de rejet du corpus (20,3~\%) est du même ordre qu'un modèle à régimes (19~\%) ou qu'un fGn de $H=0{,}6$ (16~\%).
\end{frame}

\begin{frame}{Distribution de $\hat H$ : tous les segments contre les gardés}
\footnotesize
\centering
\begin{tabular}{lccc}
\toprule
Population & $\hat H_{\mathrm{DFA}}$ & $\hat H_{\mathrm{RS}}$ & $\hat H_{\mathrm{Whittle}}$ \\
\midrule
Tous segments (145 utilisateurs) & 0,834 & 0,776 & 0,802 \\
Segments gardés (136 utilisateurs) & \textbf{0,726} & \textbf{0,723} & \textbf{0,728} \\
Gardés, pondérés par $n_{\mathrm{gardés}}$ & 0,741 & 0,731 & 0,740 \\
\bottomrule
\end{tabular}
\\[2mm]
\fig[0.50]{graphe_2_distribution.pdf}
\\[1mm]
\normalsize
Le filtre agit surtout sur la DFA ($-0{,}11$) : \alert{changement de population}, pas gain de précision.
\end{frame}

\begin{frame}{Dispersion entre utilisateurs}
\begin{columns}[T]
\column{0.52\textwidth}
\fig[0.80]{graphe_4_user_ic.pdf}
\column{0.46\textwidth}
\begin{itemize}
  \item 126 utilisateurs avec au moins 2 segments gardés ; IC à 95~\% par bootstrap sur les segments
  \item Les cinq plus gros utilisateurs portent \textbf{25,4~\%} des segments gardés
  \item Utilisateurs à faible effectif : intervalles qui se recouvrent \(\Rightarrow\) ordre non significatif
  \item $n_{\mathrm{gardés}}\ge50$ : $0{,}738\pm0{,}038$ ; $<10$ : $0{,}712\pm0{,}096$
  \item 9 utilisateurs sans segment gardé (six avec $\hat H_{\mathrm{DFA}}>1$ partout)
\end{itemize}
\end{columns}
\end{frame}

\begin{frame}{Ce que signifie $\hat H_{\mathrm{DFA}}\approx0{,}726$}
\begin{alertblock}{Message}
Ce chiffre \alert{n'est pas} une mesure de la mémoire longue du déplacement : il décrit ce que la DFA produit sur la série $d_i$ filtrée.
\end{alertblock}
\begin{itemize}
  \item Le filtre dépend de $H$ lui-même (circularité)
  \item La série mélange vitesse et rythme d'enregistrement
  \item Les segments sont fortement non stationnaires
  \item À $\rho_1$ contrôlé, la DFA distingue un fGn d'un AR(1) d'environ $0{,}07$ : du même ordre que sa RMSE \(\Rightarrow\) séparation statistique (19\,150 segments), pas pratique (segment par segment)
  \item Test AR($p$) : 20,3~\% de rejets contre 5~\% nominal, mais compatibles avec des régimes
\end{itemize}
\end{frame}

% =====================================================================
\section{Résultats sur les lois de puissance}
% =====================================================================

\begin{frame}{Résultats agrégés (137 utilisateurs)}
\footnotesize
\begin{tabular}{@{}lrrcll@{}}
\toprule
Variable & $n$ & $\hat\alpha$ & $p_{\mathrm{boot}}$ & Verdict & Meilleur AICc \\
\midrule
\code{sejour} & 5\,354 & $3{,}186\pm0{,}031$ & 0,159 & alternative meilleure & tronquée (0,967) \\
\code{distance} (même seg.) & 1\,283 & $2{,}026\pm0{,}084$ & 0,055 & rejetée & loi pure (0,528) \\
\code{sejour\_avec\_censurees} & 13\,378 & $2{,}365\pm0{,}029$ & $<0{,}001$ & rejetée & loi pure (0,577) \\
\code{distance\_a\_cheval} & 10\,216 & $1{,}779\pm0{,}012$ & $<0{,}001$ & rejetée & tronquée (0,511) \\
\bottomrule
\end{tabular}
\normalsize
\bigskip
\begin{itemize}
  \item Loi pure rejetée pour \alert{3 variables sur 4} (seul \code{sejour} n'est pas rejeté)
  \item Tronquée préférée pour \code{sejour} (poids 0,967 ; $p_{\mathrm{RV}}=0{,}0002$) et, de façon marginale, pour \code{distance\_a\_cheval} ($p_{\mathrm{RV}}=0{,}099$)
  \item Jeu de seuils : 200~m, 20~min
\end{itemize}
\end{frame}

\begin{frame}{Poids d'Akaike des cinq modèles}
\centering
\fig[0.62]{graphe_10_aicc.pdf}
\\[1mm]
\small
Aucun modèle ne s'impose pour \code{distance} et \code{sejour\_avec\_censurees} (poids $\approx0{,}5$) ; la tronquée domine pour \code{sejour}.
\end{frame}

\begin{frame}{La loi gagnante est-elle adéquate ?}
\begin{columns}[T]
\column{0.55\textwidth}
\fig[0.72]{graphe_11_pvalues.pdf}
\column{0.43\textwidth}
\small
\begin{itemize}
  \item Loi pure (Clauset) : seul \code{sejour} n'est pas rejeté ($p=0{,}159$)
  \item Tronquée (bootstrap, $B=200$) :
  \begin{itemize}
    \item \code{sejour} : \alert{rejetée} ($p<0{,}005$), rejet net
    \item \code{distance\_a\_cheval} : $p=1{,}000$, artefact ($\lambda\to0$)
    \item \code{distance} et \code{sejour\_avec\_censurees} : non rejetée, sans conclusion
  \end{itemize}
  \item Les deux tests diffèrent : leurs $p$-valeurs ne se comparent pas directement
\end{itemize}
\end{columns}
\end{frame}

\begin{frame}{Durées de séjour : coupure vers 4,7 heures}
\begin{columns}[T]
\column{0.60\textwidth}
\fig[0.60]{graphe_9_ccdf_sejour.pdf}     % <-- MODIFIÉ : 0.74 -> 0.60
\column{0.38\textwidth}
\footnotesize                              % <-- MODIFIÉ : \small -> \footnotesize
\begin{itemize}
  \item $\hat\alpha=3{,}186$ ; tronquée : $\lambda=5{,}92\times10^{-5}$~s$^{-1}$, soit $1/\lambda\approx4{,}7$~h
  \item Échelle \alert{proche du maximum observé} ($\approx5{,}5$~h) \(\Rightarrow\) interprétation ambiguë : coupure réelle, fenêtre d'observation finie, ou mélange d'utilisateurs hétérogènes
  \item La tronquée gagne en AICc (elle suit la courbure vers $10^4$~s) mais échoue à l'adéquation (les observations plongent plus vite)
  \item $\hat\alpha=3{,}19$ nettement plus grand que $\approx1{,}8$ rapporté par Song et al. : définitions et données différentes
\end{itemize}
\end{columns}
\end{frame}

\begin{frame}{Les trois autres variables}
\centering
\fig[0.40]{graphe_13_ccdf_autres.pdf}
\footnotesize
\begin{columns}[T]
\column{0.31\textwidth}
\textbf{\code{distance} (même segment)}\\
$\hat\alpha=2{,}026$, $p=0{,}055$ : rejetée ; la loi pure reste la mieux classée. Cas \alert{ambigu}.
\column{0.31\textwidth}
\textbf{\code{sejour\_avec\_censurees}}\\
$\hat\alpha=2{,}365$, $x_{\min}\approx8{,}3$~h : inclut les pauses entre segments (appareil éteint) \(\Rightarrow\) distribution probablement \alert{bimodale}.
\column{0.31\textwidth}
\textbf{\code{distance\_a\_cheval}}\\
$\hat\alpha=1{,}779$ ; $\lambda\to0$ ($1/\lambda\approx1{,}5\times10^5$~km) : la tronquée ne coupe rien. Plateau puis coupure brutale vers $10^7$~m.
\end{columns}
\end{frame}

\begin{frame}{Sensibilité aux seuils d'arrêt}
\begin{columns}[T]
\column{0.52\textwidth}
\fig[0.70]{graphe_12_sensibilite.pdf}
\column{0.46\textwidth}
\footnotesize
\begin{tabular}{@{}lrrrr@{}}
\toprule
Jeu & $n_{\mathrm{séj.}}$ & $\alpha_{\mathrm{séj.}}$ & $n_{\mathrm{dist.}}$ & $\alpha_{\mathrm{dist.}}$ \\
\midrule
100 m / 5 min & 32\,110 & 3,12 & 13\,806 & 1,89 \\
100 m / 10 min & 14\,185 & 3,18 & 4\,064 & 1,82 \\
200 m / 20 min & 5\,354 & 3,19 & 1\,283 & 2,03 \\
200 m / 30 min & 2\,789 & 3,42 & 428 & 1,54 \\
\bottomrule
\end{tabular}
\normalsize
\medskip
\begin{itemize}
  \item $\alpha_{\mathrm{séjour}}>3$ aux quatre jeux : résultat \alert{robuste}
  \item $\alpha_{\mathrm{distance}}=1{,}54$ : petit échantillon ($n=428$)
\end{itemize}
\end{columns}
\end{frame}

\begin{frame}{Lien avec l'alternance mouvement / arrêt (ON/OFF)}
\begin{block}{Relation théorique}
Phases ON/OFF à queue de survie $t^{-a}$, $1<a<2$ :
\[
H=\frac{3-a}{2}=\frac{4-\alpha}{2},\qquad 2<\alpha<3
\]
\end{block}
\small
Résultat de Taqqu, Willinger et Sherman : \emph{superposition} de nombreuses sources ; pour un individu, hypothèse de travail.
\medskip
\begin{tabular}{@{}lrlll@{}}
\toprule
Variable de durée & $\alpha$ & Dans $]2,3[$ ? & $H_{\mathrm{pred}}$ & Fiabilité \\
\midrule
\code{sejour} & 3,19 & non & --- & --- \\
\code{sejour\_avec\_censurees} & 2,37 & oui & 0,817 & très faible (bimodale) \\
\bottomrule
\end{tabular}
\medskip
\begin{itemize}
  \item \alert{Aucune prédiction solide} ; le script ON/OFF a un défaut (un seul exposant par lot de durées) : partie réelle à corriger puis régénérer
\end{itemize}
\end{frame}

% =====================================================================
\section{Discussion et conclusion}
% =====================================================================

\begin{frame}{Limites d'ensemble}
\begin{itemize}
  \item \alert{Circularité} du filtre : $\hat H_{\mathrm{gardés}}$ conditionnel, pas populationnel
  \item \alert{Validité de construit} : $d_i$ mêle vitesse et intervalle de temps
  \item \alert{Mémoire longue ou régimes} : le test de dépassement ne les distingue pas
  \item \alert{Non-stationnarité} : $\hat H$ effectif, pas celui d'un processus stationnaire
  \item \alert{Un seul jeu de données} exploitable (T-Drive écarté)
  \item \alert{Lois de puissance} : données agrégées ; test d'adéquation de la tronquée peu exigeant ; 37 utilisateurs exclus par le seuil de 300 points ; 15 utilisateurs seulement avec $\ge100$ séjours (typologie impossible en l'état)
\end{itemize}
\end{frame}

\begin{frame}{Perspectives}
\begin{block}{Départager mémoire longue et régimes (priorité)}
\small
\begin{itemize}
  \item Test de Berkes et al. : nulle = mémoire courte avec ruptures de moyenne
  \item Test de Qu : nulle = mémoire longue stationnaire contre régime ou tendance
  \item Invariance de $d$ par agrégation temporelle (Ohanissian et al.)
  \item Mesurer taille et puissance de ces tests par simulation ; modèle nul à régimes
\end{itemize}
\end{block}
\begin{columns}[T]
\column{0.48\textwidth}
\small
\textbf{Changer de série}
\begin{itemize}
  \item pas régulier (temps physique)
  \item temps d'événement (distances entre arrêts)
\end{itemize}
\column{0.48\textwidth}
\small
\textbf{Autres pistes}
\begin{itemize}
  \item segmentation par labels de transport
  \item vraisemblance avec censure
  \item typologie et cartographie des zones de mémoire
\end{itemize}
\end{columns}
\end{frame}

\begin{frame}{Conclusion}
\begin{block}{Exposant de Hurst}
$\hat H_{\mathrm{DFA}}\approx0{,}726$ sur les segments gardés de GeoLife : une persistance locale dont la robustesse aux régimes et au bruit d'échantillonnage n'est pas établie. La mémoire longue \alert{n'est pas démontrée}, ni écartée.
\end{block}
\begin{block}{Lois de puissance}
Loi pure rejetée pour 3 variables sur 4 ; la tronquée est préférée pour les séjours (coupure $\approx4{,}7$~h) mais rejetée à l'adéquation. Lien ON/OFF non testable de façon fiable.
\end{block}
\begin{block}{Apport}
Une chaîne complète, validée par simulation, qui \alert{ne cache pas ce qui limite l'interprétation}.
\end{block}
\end{frame}

\begin{frame}[plain]
\centering
\vfill
{\Huge Merci de votre attention}\\[8mm]
{\Large Questions ?}
\vfill
\end{frame}

% =====================================================================
%  ANNEXES (diapositives de secours pour les questions)
% =====================================================================
\appendix

\begin{frame}{Annexe : distribution des intervalles $\Delta t$}
\centering
\fig[0.62]{graphe_8_deltas.pdf}
\\[1mm]\small
Pics à 1--10~s : résolution des horodatages ; au-delà, décroissance monotone sans mode secondaire.
\end{frame}

\begin{frame}{Annexe : longueurs des segments retenus}
\centering
\fig[0.62]{graphe_7_longueurs.pdf}
\\[1mm]\small
Médiane de 721 points ; 66~\% des segments ont moins de 1\,000 points.
\end{frame}

\begin{frame}{Annexe : biais et RMSE des estimateurs sur fGn}
\footnotesize
\centering
\begin{tabular}{rrrrrrrr}
\toprule
$H_0$ & $N$ & \multicolumn{2}{c}{DFA} & \multicolumn{2}{c}{R/S corrigé} & \multicolumn{2}{c}{Whittle} \\
 & & biais & RMSE & biais & RMSE & biais & RMSE \\
\midrule
0,5 & 300  & $+0{,}032$ & 0,065 & $+0{,}004$ & 0,056 & $+0{,}002$ & 0,091 \\
0,7 & 300  & $+0{,}011$ & 0,061 & $-0{,}065$ & 0,093 & $-0{,}000$ & 0,095 \\
0,9 & 300  & $+0{,}011$ & 0,067 & $-0{,}149$ & 0,161 & $-0{,}002$ & 0,081 \\
0,7 & 1000 & $+0{,}007$ & 0,039 & $-0{,}048$ & 0,063 & $-0{,}001$ & 0,059 \\
0,9 & 1000 & $+0{,}002$ & 0,045 & $-0{,}121$ & 0,127 & $+0{,}001$ & 0,052 \\
\bottomrule
\end{tabular}
\end{frame}

\begin{frame}{Annexe : paramètres clés du pipeline}
\footnotesize
\begin{tabular}{@{}ll@{\hspace{10mm}}ll@{}}
\toprule
Paramètre & Valeur & Paramètre & Valeur \\
\midrule
Trou de segmentation & 900~s & Rayon / durée d'arrêt & 200~m / 1\,200~s \\
Points minimaux & 300 & Distance min. entre arrêts & 100~m \\
Z-score modifié & 3,5 & Queue minimale (Clauset) & 50 \\
Seuil de divergence & 0,15 & Bootstrap Clauset & 1\,000 \\
Échelles DFA & 20 ($4$ à $N/4$) & Seuil de rejet & $p\le0{,}1$ \\
Whittle : $m$ & $N^{0{,}65}$ & Surrogats AR($p$) & 99, $p_{\max}=5$ \\
Graine & 20261018 & Obs. min. par utilisateur & 100 \\
\bottomrule
\end{tabular}
\end{frame}

% --- OPTIONNEL : illustration (analyse faite hors mémoire ; supprimer si inutile)
\begin{frame}{Annexe (optionnel) : pourquoi un $H$ peut être aberrant, exemple du segment 0}
\small
\begin{itemize}
  \item Segment 0 : 1\,667 points ; $\hat H_{\mathrm{DFA}}=0{,}872$, $\hat H_{\mathrm{RS}}=0{,}681$, $\hat H_{\mathrm{Whittle}}=0{,}849$ ; divergence $0{,}191$ \(\Rightarrow\) \alert{suspect}
  \item Un épisode de 34 pas (2~\% du segment), probablement une dérive GPS, non marqué par le filtre
  \item En le neutralisant : DFA $0{,}747$, R/S $0{,}639$, Whittle $0{,}694$ ; divergence $0{,}108$ \(\Rightarrow\) segment gardé
  \item Placebo (34 pas ailleurs) : $\hat H$ inchangé
\end{itemize}
\medskip
Les logs signalent les points, pas l'épisode qui pèse sur $H$.
\end{frame}

\end{document}